\documentclass[11pt]{article}
\usepackage[margin=1in]{geometry}
\usepackage{booktabs}

\title{Full RL Implementation --- ISSf-CBF with Adaptive Parameters}
\author{}
\date{}

\begin{document}
\maketitle

\section*{Core Concept}

The agent learns parameters for an Input-to-State Safe Control Barrier Function (ISSf-CBF) that provides adaptive robustness margins for obstacle avoidance under dynamics disturbance and sensor uncertainty.

ISSf-CBF constraint:

\begin{verbatim}
Lgh @ u >= -alpha * h(x) + (||Lgh||^2 * phi) / max(h(x), eps)
\end{verbatim}

\begin{itemize}
\item alpha controls CBF decay rate --- how quickly the safety value can decrease
\item phi adds a robustness margin --- compensates for unmodeled disturbance and sensor error
\item The QP finds u\_safe = argmin $||$u - k\_nom$||^2$ subject to this constraint
\end{itemize}


\section*{What the Agent Controls}

Two architectures evolved over experiments:

\begin{tabular}{lll}
\toprule
Architecture & Action Space & Navigation Source \\
\midrule
Joint (exp 40-55) & [kx, ky, alpha, phi] (4D) & Learned jointly \\
Decoupled (exp 56-60) & [alpha, phi] (2D) & Frozen policy or A* planner \\
\bottomrule
\end{tabular}

\vspace{1em}
The decoupled architecture produces stronger alpha/phi learning because the navigation layer cannot absorb the safety signal.


\section*{CBF Formulations (Evolution)}

\subsection*{Original --- Squared Distance (exp 31-50)}

\begin{verbatim}
h(x) = ||robot_pos - obs_pos||^2 - estimated_radius^2
Lgh = 2 * (robot_pos - obs_pos)
\end{verbatim}

3 independent QP constraints (one per obstacle). $||$Lgh$||^2$ scales with distance$^2$ --- phi's effect is distance-dependent.


\subsection*{SHIELD SDF Pipeline (exp 53+, from Yang et al. IROS 2025, Eq 19-22)}

\begin{verbatim}
Step 1 (Eq 19): sdf(x) = min_i ||p - rho_i|| - R_i       (closest obstacle)
Step 2 (Eq 20): h_smooth = lambda * (1 - exp(-gamma * sdf(x)))  (smooth, concave)
Step 3 (Eq 21): h_hat uses projected distance along e_i from PREVIOUS timestep
Step 4 (Eq 22): piecewise extension for robot "behind" obstacle
\end{verbatim}

\begin{itemize}
\item Single QP constraint (closest obstacle only)
\item e\_i = unit direction from closest obstacle to robot (previous timestep makes h concave)
\item Lgh = lambda * gamma * exp(-gamma * proj\_dist) * e\_i
\item $||$Lgh$||^2$ is naturally directional --- phi's effect aligns with Theorem 2's projected covariance
\item Parameters: lambda = 1.0, gamma = 0.5
\end{itemize}

Piecewise extension (Eq 22):
\begin{itemize}
\item If (p - rho\_i)$^T$ * e\_i $\geq$ 0: robot in front of obstacle --- use concave approx
\item Else: robot behind obstacle --- linear approximation using gradient at boundary
\end{itemize}


\subsection*{Connection to Theorem 2 (Eq 23-24)}

Paper's constraint:
\begin{verbatim}
h_tilde(F + Gu + E[d]) - (lambda_max/2) * e^T * cov(d) * e >= alpha * h_k
\end{verbatim}

Our phi term serves the same role as the projected covariance penalty: since Lgh is proportional to e\_i, the term ($||$Lgh$||^2$ * phi) / h captures uncertainty only along the obstacle direction. The learned phi approximates the covariance magnitude that the paper computes from a CVAE.


\section*{Noise Model --- Two Sources of Uncertainty}

\subsection*{1. Dynamics Disturbance (wind / model error)}

\begin{verbatim}
bias_magnitude = U(0.3, 1.0)    # sampled once per episode
bias_angle = U(0, 2*pi)         # random direction
bias = bias_magnitude * [cos(bias_angle), sin(bias_angle)]

# Applied every timestep:
robot_pos += safe_u * dt + bias * dt
\end{verbatim}

Constant per episode --- avoids CLT averaging across timesteps. The phi term should learn to compensate for this unmodeled force.


\subsection*{2. Sensor Error (perception uncertainty)}

Two-level model (systematic calibration bias + per-measurement jitter):

\begin{verbatim}
systematic_bias = U(-0.6, 0.4)  # once per episode, all obstacles
jitter_i = U(-0.2, 0.2)         # per obstacle
error_i = systematic_bias + jitter_i

estimated_radius_i = max(true_radius_i + error_i, 1.0)
\end{verbatim}

\begin{itemize}
\item true\_radius: used for collision detection (ground truth)
\item estimated\_radius: what the robot believes, used in CBF computation
\item Negative error = underestimate (dangerous: CBF thinks obstacle is smaller)
\item Positive error = overestimate (conservative: robot thinks obstacle is bigger)
\item Systematic bias is coherent across all obstacles --- stronger learning signal than independent errors
\end{itemize}


\section*{Observation Space}

14-dimensional for 3 obstacles (20D for 5 obstacles):

\begin{verbatim}
obs = [
    # Per obstacle (3 or 5):
    obs_x - robot_x,           # relative x position
    obs_y - robot_y,           # relative y position
    estimated_radius,          # what robot believes (includes sensor error)
    # ... repeated for each obstacle

    # Target:
    target_x - robot_x,       # relative x to target
    target_y - robot_y,       # relative y to target
    target_radius,             # acceptance radius

    # Velocity:
    vx, vy                     # current robot velocity
]
\end{verbatim}

All positions are relative --- the policy sees distances, not absolute coordinates.


\section*{Environment Setup}

\begin{tabular}{ll}
\toprule
Parameter & Value \\
\midrule
Timestep dt & 0.1 s \\
Max speed & 6.0 m/s (action bounds [-6, 6]) \\
Target position & (150.0, U(-5, 5)) \\
Max steps per episode & 800 \\
Out-of-bounds & x in [-5, 160], y in [-20, 20] \\
Number of obstacles & 3 (or 5 in exp 60) \\
Obstacle x-bands & [20,50], [60,100], [110,140] \\
Obstacle y-range & U(-10, 10) (or U(-4, 4) tight) \\
Obstacle radius & U(3.0, 7.0) per obstacle \\
\bottomrule
\end{tabular}


\section*{Domain Randomization}

All parameters sampled per episode to ensure generalization:

\begin{tabular}{lll}
\toprule
Parameter & Range & Scope \\
\midrule
Obstacle positions & x-bands, y-range & Per obstacle \\
Obstacle true radius & U(3.0, 7.0) & Per obstacle \\
Sensor systematic bias & U(-0.6, 0.4) & Per episode \\
Sensor jitter & U(-0.2, 0.2) & Per obstacle \\
Dynamics bias magnitude & U(0.3, 1.0) & Per episode \\
Dynamics bias angle & U(0, 2*pi) & Per episode \\
Target y-position & U(-5, 5) & Per episode \\
Target radius & U(1.5, 3.0) & Per episode \\
\bottomrule
\end{tabular}


\section*{Reward Structure}

\begin{verbatim}
if collision (true_dist < 0):
    reward = -250.0
    terminated = True
else:
    progress = prev_dist2target - dist2target
    reward = progress * 50.0 - 1.0   # progress toward goal, time penalty

if reached_target:
    reward += 100.0
    terminated = True

if out_of_bounds:
    reward -= 50.0
    terminated = True
\end{verbatim}

The collision penalty magnitude affects phi learning: higher penalty (-250 vs -100) creates more pressure to use phi for safety margins.


\section*{Robot Dynamics}

Single integrator --- position is directly controlled by velocity:

\begin{verbatim}
robot_pos += safe_u * dt + bias * dt
\end{verbatim}

safe\_u comes from the QP (modified k\_nom). bias is the unmodeled dynamics disturbance.


\section*{QP Safety Filter}

\begin{verbatim}
u_safe = argmin ||u - k_nom||^2
    subject to: Lgh @ u >= -alpha * h + (||Lgh||^2 * phi) / max(h, eps)
\end{verbatim}

\begin{itemize}
\item Squared CBF: 3 constraints (one per obstacle)
\item SHIELD SDF: 1 constraint (closest obstacle only)
\item Solver: CVXPY with OSQP backend, warm-started
\item Fallback: u\_safe = [0, 0] on solver failure
\item Output clipped to [-max\_speed, max\_speed]
\end{itemize}


\section*{Navigation Strategies (Decoupled Architecture)}

\subsection*{Frozen Learned Policy (exp 56, 58)}

Phase 1: Train navigation policy: obs $\rightarrow$ [kx, ky] (with or without obstacles).

Phase 2: Freeze nav policy, train safety policy: obs $\rightarrow$ [alpha, phi].

The frozen navigator drives the robot near obstacles because it cannot adapt --- the CBF must intervene, providing consistent gradient signal for alpha/phi.

\subsection*{A* Path Planner (exp 57, 59, 60)}

Classical A* on an occupancy grid:
\begin{itemize}
\item Grid resolution: 1.0 m/cell
\item Obstacle inflation: estimated\_radius + margin (0.5-1.0 m)
\item 8-connected neighbors
\item Replan every 50 timesteps
\item Waypoint following: proportional controller with K\_p = 4.0, threshold 3.0 m
\end{itemize}

A* provides optimal shortest paths and will route through tight spaces, guaranteeing near-obstacle training signal.


\section*{Training Configuration}

\begin{tabular}{ll}
\toprule
Parameter & Value \\
\midrule
Algorithm & PPO (Stable Baselines 3) \\
Policy network & MlpPolicy \\
Device & CUDA (RTX 5090) \\
Rollout steps & 4096 per env \\
Parallel environments & 16 (SubprocVecEnv) \\
Total timesteps & 10M (5M for easy nav phase) \\
\bottomrule
\end{tabular}


\section*{Key Findings}

\begin{enumerate}
\item Alpha learns proximity-adaptive behavior when the CBF constraint activates frequently.
\item Phi stays flat when alpha is available as a competing safety knob (joint training).
\item Fixed alpha forces phi to learn --- removing alpha from the action space makes phi the only safety lever.
\item Higher collision penalty (-250 vs -100) creates more phi activity.
\item Decoupled architecture (frozen nav or A*) gives the cleanest alpha/phi signal --- the navigation layer cannot absorb the safety gradient.
\item The robot avoids rather than slows down in open environments --- environment geometry determines whether the CBF is stressed.
\item SHIELD SDF provides a theoretically grounded single-constraint CBF where phi's effect naturally aligns with the projected covariance from Theorem 2.
\end{enumerate}

\end{document}
